% ID: SEL-D0047-OB-E01
% ESTADO: cotejada. Segunda pasada visual y borrador compilado revisados.
% Fuente: Selectividad/sel_2014_matematicas_aplicadas/Matematicas_CCSS_SEPT.pdf
% SHA-256: ee5bd8e23465d55ed58339e7d9adf3d58423ac2966b53d92e39a2f9b35558ab1
% Página: 2; opción B; ejercicio original 1.

\begin{enumerate}[label=\alph*)]
\item \textbf{(1.5 puntos)} Plantee, sin resolver, el siguiente problema:

``Un mayorista vende productos congelados que presenta en envases de dos tamaños, pequeños y grandes. La capacidad de sus congeladores no le permite almacenar más de 1000 envases en total. En función de la demanda sabe que debe mantener un stock mínimo de 100 envases pequeños y 200 envases grandes. La demanda de envases grandes es igual o superior a la de envases pequeños. El coste por almacenaje es de 10 céntimos de euro por cada envase pequeño y de 20 céntimos de euro por cada envase grande. \textquestiondown{}Qué número de envases de cada tipo proporciona el mínimo coste de almacenaje?''
\item \textbf{(1 punto)} Represente el recinto que determinan las inecuaciones
\[
2x\geq10+y,\qquad x\leq2(5-y),\qquad x\geq0,\qquad y\geq0.
\]
\end{enumerate}
